\section{Evaluation Plan}
\label{sec:evaluation}

The evaluation tests each mechanism in a workload that exposes its
expected benefit, then asks whether those benefits improve fork,
migration, and concurrent execution. The formal study remains to be
conducted. Section~\ref{sec:agent-observations} reports the completed
request characterization; an exploratory replay is separately archived
in the companion experiment repository.

\subsection{Systems, inputs, and accounting}

JuiceFS is the external distributed-file-system
baseline~\cite{juicefs_architecture}. Internal comparisons use ordinary
volumes, native-metadata workspaces, and packed-metadata workspaces.
Mechanism ablations keep the execution engine and unrelated policies
fixed; comparisons between complete paths are reported separately.
Resource pools match CPU, memory, storage, and metadata-service budgets.
Concurrent comparisons require independent services; runs sharing a
service execute sequentially. Software revisions, effective layouts,
placement, and supported operations are recorded.

Inputs include coding and build/test traces, a million-file collection
with 4 KiB payloads, and fixed collections shared by concurrent readers.
The companion protocols specify generation, schedules, and repetitions.
ML dataset access, repeated publication, retention, and scale studies
extend the configured replay and small-file experiments.

Accounting covers preparation, execution, validation, retention, and
reclamation, including failed attempts. Shared preparation is counted
once per baseline; object residence continues until reclamation.
CPU time, transferred bytes, and monetary cost retain separate units.
Batch duration is measured directly, with queueing included, rather
than summed from overlapping tasks. Replay and actual agent execution
are separate studies.

\subsection{Overlay fork and retained branches}

\researchq{1} How much state and work does overlay fork save without
packing metadata?
Many children start from one revision and modify a small part of its
tree. A native-metadata workspace is compared with an eagerly
materialized private view in the same engine, preserving the API,
payload sharing, caches, and durability semantics. An ordinary copied
volume provides a separate system-level comparison.

Increasing fan-out tests preparation sharing. Star, branching, and
chained histories use equal fork budgets to expose the cost of
descendants of modified parents. Measurements include baseline capture,
creation to first correct operation, first durable edit, save latency,
private records, and uploaded bytes. Readback checks inheritance and
parent/sibling isolation. Retaining, reopening, and deleting states
measures accumulated metadata and actual reclamation; admission failures
and outstanding cleanup remain visible.

\subsection{Read-only metadata service}

\researchq{2} When does packing a fixed small-file collection reduce
database load enough to repay construction and publication?
Concurrent readers repeatedly access a fixed collection while other
workspaces update the same metadata service. The pure metadata
comparison changes native versus packed representation within one
engine, holding payload layout and caches identical. A separate
comparison permits each complete format to use its own data layout.
Standalone read-only mounts isolate the representation from upper checks.

Native version capture and binding are timed alongside the additional
construction, upload, validation, and publication required for an
unprepared packed version. First and later forks separate preparation
from reuse. Equal-use comparisons vary consumers, version churn, and
post-use retention to identify resource-specific break-even regions.
Query offload and persistent database reduction are measured separately:
retained source records still occupy space.

Positive and negative lookups, inode attributes, neighboring names,
enumeration followed by attributes, and lookup followed by reads
exercise different locality. Cold, index-warm, and metadata-warm runs
declare supported cache budgets. Counters separate lower-database,
private-upper, and control traffic from index, joint-page, descriptor,
and payload requests. Latency, database CPU and memory, persistent
bytes, transfers, decoding, and verification reveal where work moves.
A fixed-rate mutable-workspace probe measures interference through its
queue, deadline completion, and tail latency.

An ML-oriented case separates dataset discovery from sample loading:
fresh enumeration versus sample-list reuse, then ordered versus
shuffled sample access, with outputs outside the fixed dataset.
Larger-file inputs test the role of per-file metadata. Compatible
Nydus or SOCI readers provide read-only layout comparisons at matched
content, integrity, cache, and payload settings. An immutable sorted-KV
control isolates joint-page organization. These comparisons charge
construction, updates, and residence as well as query service.

\subsection{Object formation and backend work}

\researchq{3} How do client and backend granularity affect physical
work, write completion, and repair?
Small reads within larger objects expose excess fetching and
verification; scans test the request overhead of smaller units.
Identical writes compare client sizes of 256 KiB, 1 MiB, and 4 MiB at
fixed backend geometry. Fixed finalized objects and S3 requests then
compare admitted encoding-block hints. Verification policies are varied
separately at a fixed effective encoding block size. Range and
whole-object fetching have matched cache, coalescing, and prefetch
controls.

Boundary-crossing reads, shared and disjoint keys, overwrites, and
supported degraded reads test the resulting tradeoffs. Object counts,
finalized lengths, HTTP and device bytes, checksum and encoding CPU,
durable-write latency, and repair traffic explain performance.
Actual $(g,L_o,B,q)$ is recorded. An unchanged $(B,q)$ provides no
backend-layout intervention, even if the policy changes client fetching.
HTTP bytes do not substitute for missing physical-read counters.

\subsection{Migration and recovery}

\researchq{4} How much work does migration require before and after
publication?
An agent resumes on another worker while neighboring tasks continue.
Only paths that reopen the existing workspace enter this comparison;
fresh-only paths are evaluated for creation instead. Paired stay/move
runs vary the active population and migrated fraction. Source and target
host identities establish that a move crosses workers; JuiceFS uses its
native directory reopen operation.

Handoff occurs at a completed-tool boundary, with either a published
version or pending modifications. Interruption includes queueing,
draining, capture, any required packing, publication, and authority
transfer. Cold and warm targets measure first correct operation,
continuation cost, transfers, and neighboring tasks' completion and
tail latency.

File-state tests check names and bytes. Actual agent tests preserve
session history and remaining budget through export/import, require
a subsequent valid tool result, and grade final task completion.
Published file state does not checkpoint process memory or open handles.
Failures before and after publication, disconnected source coordination,
and target failure test recovery and stale-writer rejection.
Independent readback distinguishes storage-enforced rejection from
process termination or application checks.

\subsection{Concurrent tasks and analysis}

\researchq{5} Do component savings increase useful capacity at fixed
resources?
A fixed assignment set runs at increasing concurrency under a common
deadline, including queueing. A separate scale study increases workers
and assigned work against fixed backend capacity. Placement and
replay-engine CPU limits are recorded.

Internal ablations change one supported mechanism; composition is
compared with its best supported proper subset. External system
comparisons remain separate. Paired blocks report raw outcomes and
ranges. Stronger inference uses a sample size chosen from a separate
variance pilot; requests and sibling tasks are not independent
replicates. Deadline completion includes failures in its denominator.
Complete-success makespan is reported only when every assignment
succeeds, alongside failed outcomes and measurement gaps.
